\chapter{Background}
\label{cha:1}
I will intruduce concepts like foas, hoas, phoas and verification condition generation. Then I will motivate using phoas for verification condition generation

\section{Representing abstract syntax}

\subsection{First Order Abstract Syntax}
Here I will describe the first order approach for representing abstract syntax, I can describe nominal string representations/de bruijn levels and mention their strenghts and weaknesses. (information from eg. https://jesper.sikanda.be/posts/1001-syntax-representations.html)
\subsection{Higher order abstract syntax}
Describe HOAS representation 



\subsection{Parametric Higher order abstract syntax}

(information from eg. Parametric Higher-Order Abstract Syntax for Mechanized
Semantics)

\subsection{Translations between representations}
Describe translations phoas<->foas and phoas-> hoas such as described in Unembedding Domain-Specific Languages.


\section{monads?}
(I'm not sure if this is necessary but it is something most other students would not have as background knowledge)
Describe option, reader, state and continuation monads.

\section{Verification condition generation}
Describe what a verification condition is and how it can be computed by recursively calculating the weakest precondition.

\subsection{Wstore and its combinators}
Describe the Wstore monad and assert, demonic and angelic choice combinators.
\subsection{generating verification conditions}
Describe how to use these combinators to generate the weakest precondition.



(Verified Symbolic Execution with Kripke Specification
Monads (and No Meta-programming)


